
\section{Introduction}

Neutrino- and antineutrino-induced interactions at energies of
a few GeV are a proving ground for weak interaction
phenomenology in nuclei~\cite{Alvarez-Ruso:2014bla}.
Measurements in this energy range are important because 
neutrino oscillation experiments~\cite{Ayres:2004js,LBNE} need 
detailed understanding of the large variety of processes allowed.  
As a result, there is a growing body of new high-quality measurements for 
neutrino-nucleus
interactions~\cite{AguilarArevalo:2007ab,Lyubushkin:2008pe,nubarprl,miniboone_piprod}. 
In particular, neutrino charged-current neutral pion 
production has become an important benchmark providing new challenges 
to theories describing this process~\cite{gibuu-pi,valencia-pi,Rodrigues:2014jfa,Yu:2014yja}.

Most of the published data for pion production in nuclei uses 
neutrino beams.  Neutral pion production in nuclei for anti-neutrinos,
however, is much less studied.
Only one data point for this channel in the few-GeV energy range exists 
in the literature: 
a measurement of $\bar{\nu}_\mu p \rightarrow \mu^+ n \pi^0$
from SKAT in a heavy liquid (CF$_3$Br) bubble chamber based on 20 events
at an average neutrino energy of 7 GeV~\cite{Grabosch89}. 
Measurements of $\pi^0$ production by neutrinos have been made in 
deuterium bubble chambers~\cite{Allasia90, Barish79, Radecky82, Kitagaki86}
 and more recently on nuclear targets in the 0.1 - 1 GeV energy range
using 
the MiniBooNE detector~\cite{AguilarArevalo:2008qa} with a mineral oil (CH$_2$) target, and
using the SciBar detector in K2K and SciBooNE 
experiments~\cite{Mariani:2007,Kurimoto:2009} with plastic scintillator (CH).
These recent measurements, as well as data on charged pion 
production~\cite{miniboone_piprod,Eberly:2014mra} and 
neutral pion production~\cite{Augil11} have been difficult 
for event generators and theoretical calculations to 
describe accurately~\cite{gibuu-pi,valencia-pi,Rodrigues:2014jfa,Yu:2014yja}.
%establishing the need for further studies.
Every prediction must have a model for $\pi^0$ production from
nucleons; all use isospin decompositions within
a helicity formalism that are tuned to available data~\cite{Rein:1980wg}.

Charged-current single $\pi^0$ ($1\pi^0$) production in the few-GeV region
is modeled both as decays of nucleon resonances (most strongly
the $\Delta(1232)$) and nonresonant processes. The production 
in nuclei can be
either direct, through the reaction $\bar{\nu}_\mu p \rightarrow \mu^{+} n
\pi^0$, or indirect, for example, through charge exchange (CEX) of a charged pion in the nucleus, 
%$p(\pi^- , \pi^{0})n$ or $n(\pi^+, \pi^{0})p$. 
$p\pi^-\rightarrow n\pi^0$ or $n\pi^+\rightarrow p\pi^0$.
%Although these are reasonable assumptions, they are not proven. 
New data will provide useful tests  
of both neutrino-induced resonance production and final state 
interaction (FSI) models. 

Neutrino interaction measurements are also important to the analysis of neutrino
oscillation experiments~\cite{Abe:2011ks,Ayres:2004js,AguilarArevalo:2010wv,LBNE}. 
These experiments require the neutrino flavor be
identified and the neutrino energy to be reconstructed on an event-by-event basis.
An accurate modeling of these particles requires knowledge of both the underlying
neutrino-nucleon interactions and of the final-state modifications that arise within the
target nuclei (such as carbon) of which the massive oscillation detectors are comprised.
%Neutrino interaction measurements are important components of the
%neutrino oscillation experiments analyses~\cite{Abe:2011ks,Ayres:2004js,AguilarArevalo:2010wv,LBNE}, which
%must identify neutrino flavor and reconstruct 
%neutrino energy based on the final state particles observed in the detector. 
%Correctly modeling the identity and kinematics of these particles requires 
%accurate knowledge of both neutrino-nucleon interaction processes and the 
%modifications that result when the interaction takes place in a nucleus, 
%as is the case for current and future oscillation experiments. 
Pion production is a source of backgrounds and systematic uncertainties in neutrino oscillation experiments. 
Neutral-current $\pi^0$ production, for example, is a dominant background in
$\nu_e$ ($\bar{\nu}_e$) appearance experiments because the $\pi^0$ can mimic 
a final state electron (positron).
In addition, experiments that reconstruct neutrino energy by 
identifying quasielastic events,
~$\nu_\ell N(n) \rightarrow \ell^{-} p$~
or~$\bar{\nu}_\ell N(p) \rightarrow \ell^{+} n$, 
interactions in which a pion is produced but then absorbed in the target 
nucleus can be mistaken for quasielastic signal and yield an incorrect estimate of the 
incident neutrino energy.

New measurements of $1\pi^0$ production by 
charged-current $\bar{\nu}_\mu$ interactions in plastic scintillator (CH) using
the \minerva detector are presented.
Flux-integrated single differential cross sections as a function of $\pi^0$ momentum and production angle for events with a
single observed $\pi^0$ and no $\pi^\pm$ exiting the interaction nucleus have
been measured and are compared to
predictions from the GENIE~\cite{Andreopoulos201087}, NuWro~\cite{Juszczak:2005zs,Golan:2012rfa}, 
and NEUT~\cite{Hayato:2009zz} event generators.
